% ============================================================
% Section 66: 自由电子模型——费米球与布里渊区边界的接触
% ============================================================
\newpage
\section{固体理论：自由电子模型——费米球与布里渊区边界的接触}
\label{sec:fermi-brillouin-contact}

\begin{modelbox}[题目：Cu-Zn合金中费米球与布里渊区边界接触的条件]
向铜中掺锌，一些铜原子被锌原子所取代。采用自由电子模型，求锌原子与铜原子之比为何值时，费米球与第一布里渊区边界相接触？（铜是面心立方晶格，单价，锌是二价。）
\end{modelbox}

\begin{tipbox}[考点快速分析]
\begin{itemize}
    \item \textbf{自由电子模型}：价电子视为自由电子气，费米波矢 $k_F$ 由电子数密度 $n$ 决定
    \item \textbf{布里渊区边界}：FCC的第一布里渊区是截角八面体，沿$\langle 111\rangle$方向边界最近
    \item \textbf{合金电子浓度}：Zn替代Cu后，平均每个原子贡献的价电子数增加，费米球膨胀
    \item \textbf{Hume-Rothery规则}：当费米球接触布里渊区边界时，合金的相稳定性发生变化
\end{itemize}
\end{tipbox}

\begin{derivationbox}[标准解题过程]
\textbf{Step 1: 电子浓度的计算}

设晶体中共有 $M$ 个格点。掺入锌的原子百分比为 $x$，则铜的原子百分比为 $1-x$。
\begin{itemize}
    \item 锌为二价，每个锌原子贡献 $2$ 个自由电子
    \item 铜为单价，每个铜原子贡献 $1$ 个自由电子
\end{itemize}
晶体中总自由电子数为
\begin{equation}
N = 2(xM) + 1\cdot(1-x)M = (1+x)M.
\end{equation}

面心立方晶格每个惯用晶胞含 $4$ 个格点，晶胞体积为 $a^{3}$。因此每个格点平均占据的体积为 $a^{3}/4$，晶体总体积为 $V=M\cdot a^{3}/4$。

自由电子数密度为
\begin{equation}
n = \frac{N}{V} = \frac{(1+x)M}{M\cdot a^{3}/4} = \frac{4(1+x)}{a^{3}}.
\label{eq:electron-density}
\end{equation}

\textbf{Step 2: 费米波矢}

在自由电子模型中，考虑自旋简并后，电子数密度与费米波矢的关系为
\begin{equation}
n = \frac{k_F^{3}}{3\pi^{2}},
\qquad\text{即}\qquad
k_F = (3\pi^{2}n)^{1/3}.
\end{equation}

\textbf{Step 3: 第一布里渊区边界的位置}

面心立方晶格的倒格子是体心立方，倒格常数为 $4\pi/a$。第一布里渊区是倒空间中的截角八面体。

费米球与布里渊区边界接触，发生在费米球半径等于布里渊区中心到边界的最短距离。对于面心立方晶格，该最短距离沿 $\langle 111\rangle$ 方向，等于倒格子原胞体心到六边形面中心的距离：
\begin{equation}
k_{\mathrm{边界}} = \frac{1}{4}\sqrt{3}\left(\frac{4\pi}{a}\right) = \frac{\sqrt{3}\,\pi}{a}.
\end{equation}
（推导：BCC倒格子中，沿$\langle 111\rangle$方向的倒格矢长度为 $\sqrt{3}\cdot 4\pi/a$，布里渊区边界位于该矢量的中点，故距离为一半。）

\textbf{Step 4: 求解掺杂比例}

令费米球半径等于布里渊区边界距离：
\begin{equation}
k_F = k_{\mathrm{边界}}
\quad\Longrightarrow\quad
(3\pi^{2}n)^{1/3} = \frac{\sqrt{3}\,\pi}{a}.
\end{equation}
两边立方：
\begin{equation}
3\pi^{2}n = \frac{3\sqrt{3}\,\pi^{3}}{a^{3}}
\quad\Longrightarrow\quad
n = \frac{\sqrt{3}\,\pi}{a^{3}}.
\label{eq:critical-density}
\end{equation}

将电子密度表达式 \eqref{eq:electron-density} 代入 \eqref{eq:critical-density}：
\begin{equation}
\frac{4(1+x)}{a^{3}} = \frac{\sqrt{3}\,\pi}{a^{3}}.
\end{equation}
消去 $a^{3}$，解得
\begin{equation}
1+x = \frac{\sqrt{3}\,\pi}{4}
\approx \frac{1.732\times 3.142}{4}
\approx 1.360,
\end{equation}
即
\begin{equation}
\boxed{x \approx 0.36 = 36.0\%.}
\end{equation}

锌与铜的原子数之比为
\begin{equation}
\frac{x}{1-x} = \frac{0.36}{0.64} = \frac{9}{16}.
\end{equation}
\textit{结论}：当锌的原子百分比约为 $36\%$（锌铜原子数之比为 $9:16$）时，费米球恰好与第一布里渊区边界接触。
\end{derivationbox}

\begin{formulabox}[答案汇总]
\begin{center}
\begin{tabular}{ll}
\toprule
\textbf{问题} & \textbf{答案} \\
\midrule
锌的原子百分比 & $x \approx 36\%$ \\
锌与铜的原子数之比 & $\mathrm{Zn}:\mathrm{Cu} = 9:16$ \\
\bottomrule
\end{tabular}
\end{center}
\end{formulabox}

\begin{insightbox}[物理图像：为什么接触布里渊区边界如此重要？}
当费米球膨胀到与布里渊区边界接触时，电子结构发生质变：
\begin{itemize}
    \item \textbf{能隙打开}：布里渊区边界处存在能带间隙，费米面附近的电子态密度出现突变
    \item \textbf{费米面畸变}：球形的费米面在边界处被``压扁''，偏离自由电子的抛物线行为
    \item \textbf{物理性质异常}：电导率、磁化率、比热等物理量出现非单调变化
\end{itemize}

\textbf{Hume-Rothery电子浓度规则}：
许多合金系统的相稳定性与电子浓度（每个原子的平均价电子数）密切相关：
\begin{itemize}
    \item FCC $\alpha$相（如Cu基固溶体）：电子浓度上限约 $1.36\,e^-/\text{atom}$
    \item BCC $\beta$相：电子浓度约 $1.48\,e^-/\text{atom}$
    \item 复杂立方 $\gamma$相：电子浓度约 $1.54\,e^-/\text{atom}$
\end{itemize}

本题中，Cu-Zn合金在 $x=36\%$ 时平均价电子数为 $1+x=1.36$，恰好对应$\alpha$相的溶解度极限。实际Cu-Zn合金中$\alpha$相的最大锌溶解度约为 $38\%\sim 39\%$，与理论值非常接近——这是自由电子模型在合金理论中最成功的应用之一。
\end{insightbox}

\begin{thinkbox}[考场速记：自由电子+布里渊区问题的通用模板]
遇到``费米球接触布里渊区边界''类问题：
\begin{enumerate}
    \item \textbf{写电子浓度}：$n = $(平均价电子数)$/v_{\text{原子}}$
    \item \textbf{写费米波矢}：$k_F = (3\pi^{2}n)^{1/3}$
    \item \textbf{写布里渊区边界}：
    \begin{itemize}
        \item FCC：$k_{\text{边界}} = \sqrt{3}\,\pi/a$（沿$\langle 111\rangle$）
        \item BCC：$k_{\text{边界}} = \sqrt{2}\,\pi/a$（沿$\langle 110\rangle$）
        \item sc：$k_{\text{边界}} = \pi/a$（沿$\langle 100\rangle$）
    \end{itemize}
    \item \textbf{令 $k_F = k_{\text{边界}}$}，解出临界电子浓度
    \item \textbf{换算合金比例}
\end{enumerate}

\textit{布里渊区边界速查}：
\begin{center}
\begin{tabular}{lcc}
\toprule
\textbf{结构} & \textbf{边界方向} & \textbf{边界距离} \\
\midrule
简立方（sc） & $\langle 100\rangle$ & $\pi/a$ \\
面心立方（fcc） & $\langle 111\rangle$ & $\sqrt{3}\,\pi/a$ \\
体心立方（bcc） & $\langle 110\rangle$ & $\sqrt{2}\,\pi/a$ \\
\bottomrule
\end{tabular}
\end{center}
\end{thinkbox}
